\documentclass[border=4pt]{standalone}
\usepackage[siunitx]{circuitikz}
\usetikzlibrary{calc}

\begin{document}
\begin{circuitikz}[font=\small,line width=0.9pt]
  \ctikzset{tripoles/thickness=1,logic ports/thickness=1,
    flipflops/thickness=1}
  \node[font=\bfseries] at (6.0,6.5)
    {Asynchronous-assert, synchronous-deassert reset-release cell};

  \node[flipflop D,flipflop def={td={CLR},nd=1},scale=1.05] (F1) at (3.8,3.7) {};
  \node[flipflop D,flipflop def={td={CLR},nd=1},scale=1.05] (F2) at (7.8,3.7) {};
  \node[above] at (F1.north) {release stage 1};
  \node[above] at (F2.north) {release stage 2};

  \coordinate (ONE) at ($(F1.pin 1)+(-2.2,0)$);
  \draw (ONE) node[left] {logic 1} -- (F1.pin 1);
  \draw (F1.pin 6) -- (F2.pin 1) node[midway,above] {$Q_1$};
  \draw (F2.pin 6) -- ++(1.4,0) node[right] {$RESET_{domain\_n}=Q_2$};

  \draw (1.1,0.8) node[left] {domain $CLK$} -- (8.8,0.8);
  \draw (2.8,0.8) |- (F1.pin 3);
  \draw (6.8,0.8) |- (F2.pin 3);
  \fill (2.8,0.8) circle (1.4pt);
  \fill (6.8,0.8) circle (1.4pt);

  \draw (F1.down) -- ++(0,-0.65) node[below] {$RESET_n$};
  \draw (F2.down) -- ++(0,-0.65) node[below] {$RESET_n$};

  \node[align=center] at (6.0,-1.05)
    {Matching $RESET_n$ labels denote one external active-low reset net.\\
     $RESET_n=0$ clears both stages after clear-to-Q delay, without a clock.\\
     Normally logic 1 advances on the next two rising edges; close release adds variable latency.};
\end{circuitikz}
\end{document}
